\chapter[Работа с библиотеками Soapy SDR, Libiio.]{Работа с библиотеками Soapy SDR, Libiio.\\ Формирование и передача с SDR сигналов\\ произвольной формы}

\begin{leftbar}
    \href{https://github.com/someotb/sdr}{Ссылка на GitHub}
\end{leftbar}

\section*{Цель практики:}

Лучше разобраться с библиотеками, сформировать сэмплы и отправить их, визулизировать полученые данные.

\section*{Краткие теоретические сведения}

Вспомним структуру отправляемого буфера

\begin{figure}[H]
    \centering
    \includegraphics[width=1.0\textwidth]{buffer_structer.png}
    \caption{Структура буфера}
\end{figure}

Структура сэмплов

\begin{figure}[H]
    \centering
    \includegraphics[width=1.0\textwidth]{samples_adalmpluto.png}
    \caption{Структура сэмплов}
\end{figure}

Как читаются сэмплы

\begin{figure}[H]
    \centering
    \includegraphics[width=1.0\textwidth]{read_sample.png}
    \caption{Чтение сэмплов}
\end{figure}

\section*{Выполнение}

Сперва я заполнил буферы значениями 1500 << 4 ($1500 * 2^4 = 24000$)
Передавал их на SDR и смотрел что получается.
Визуализировать удобно с помощью кода на python

\begin{lstlisting}[language=python, caption=Код для визуализации полученных данных с SDR]
import matplotlib.pyplot as plt
import numpy as np

rx = np.fromfile("/home/plutoSDR/sdr/pluto/dev/rx.pcm", dtype=np.int16)

samples = []
for x in range(0, len(rx), 2):
    samples.append(rx[x] + 1j * rx[x+1])

samples = np.array(samples)

ampl = np.abs(samples)
phase = np.angle(samples)
time = np.arange(len(samples))

plt.figure(figsize=(10, 8))

plt.subplot(3, 1, 1)
plt.plot(time, ampl)
plt.title("Amplitude")
plt.ylabel("Amplitude")
plt.grid(True)

plt.subplot(3, 1, 2)
plt.plot(time, phase)
plt.title("Phase")
plt.ylabel("Phase (rad)")
plt.grid(True)

plt.subplot(3, 1, 3)
plt.plot(time, np.real(samples), label='I (Real)')
plt.plot(time, np.imag(samples), label='Q (Imag)')
plt.title("I and Q components")
plt.ylabel("Value")
plt.xlabel("Time (samples)")
plt.legend()
plt.grid(True)

plt.tight_layout()
plt.show()
\end{lstlisting}

На графике получаем, амплитуду, фазу и I,Q части

\begin{figure}[H]
    \centering
    \includegraphics[width=0.9\textwidth]{results.png}
    \caption{Графики}
\end{figure}


\section*{Вывод}
Самостоятельно записали произвольные данные в буферы, который потом отправили на SDR, получили их, вывели на графики, с помощью
python кода, наглядно показано, что получаем на выходе после всех этапов отправки данных.